\documentclass[10pt, letterpaper]{article}

\usepackage[
    ignoreheadfoot,
    top=1.45cm,
    bottom=1.45cm,
    left=1.65cm,
    right=1.65cm,
    footskip=0.8cm
]{geometry}

\usepackage{titlesec}
\usepackage{tabularx}
\usepackage{array}
\usepackage[dvipsnames]{xcolor}
\definecolor{primaryColor}{RGB}{0,0,0}
\usepackage{enumitem}
\usepackage{amsmath}
\usepackage[
    pdftitle={Jamie Wei Resume},
    pdfauthor={Jamie Wei},
    pdfcreator={LaTeX},
    colorlinks=true,
    urlcolor=primaryColor
]{hyperref}
\usepackage{bookmark}
\usepackage{changepage}
\usepackage{needspace}
\usepackage{iftex}

\ifPDFTeX
    \input{glyphtounicode}
    \pdfgentounicode=1
    \usepackage[T1]{fontenc}
    \usepackage[utf8]{inputenc}
\fi

\usepackage{charter}

\raggedright
\pagestyle{empty}
\setcounter{secnumdepth}{0}
\setlength{\parindent}{0pt}
\setlength{\topskip}{0pt}
\pagenumbering{gobble}

\titleformat{\section}
    {\needspace{3\baselineskip}\bfseries\large}
    {}
    {0pt}
    {}
    [\vspace{1pt}\titlerule]

\titlespacing{\section}{0pt}{0.20cm}{0.12cm}

\renewcommand\labelitemi{$\vcenter{\hbox{\small$\bullet$}}$}

\newenvironment{highlights}{
    \begin{itemize}[
        topsep=0.04cm,
        parsep=0.02cm,
        partopsep=0pt,
        itemsep=0.03cm,
        leftmargin=12pt
    ]
}{
    \end{itemize}
}

\newenvironment{onecolentry}{
    \begin{adjustwidth}{0cm}{0cm}
}{
    \end{adjustwidth}
}

\begin{document}

%%%%%%%%%%%%%%%%%%%%%%%%%%%%%%%%%%%%%
% HEADER
%%%%%%%%%%%%%%%%%%%%%%%%%%%%%%%%%%%%%

\begin{center}
    {\Huge\textbf{Jamie Wei}}\\[4pt]
    U.S. Citizen
    \quad|\quad
    \href{tel:+16787620657}{(678) 762-0657}
    \quad|\quad
    \href{mailto:Jamiejwei@gmail.com}{Jamiejwei@gmail.com}\\[3pt]
    \href{https://www.linkedin.com/in/jamie-wei-5a4a31378/}
        {linkedin.com/in/jamie-wei-5a4a31378}
    \quad|\quad
    \href{https://github.com/JamieWei213213}
        {github.com/JamieWei213213}
\end{center}

%%%%%%%%%%%%%%%%%%%%%%%%%%%%%%%%%%%%%
% EDUCATION
%%%%%%%%%%%%%%%%%%%%%%%%%%%%%%%%%%%%%

\section{Education}

\begin{tabularx}{\textwidth}{Xr}
    \textbf{University of Southern California}
        & \textit{Jan 2025 -- Dec 2026 (Expected)}\\
    Master of Science in Applied Data Science & \\
\end{tabularx}

\vspace{0.08cm}

\begin{tabularx}{\textwidth}{Xr}
    \textbf{University of California San Diego}
        & \textit{Sep 2019 -- Mar 2023}\\
    Bachelor of Science in Cognitive Science,
    Machine Learning Specialization & \\
\end{tabularx}

%%%%%%%%%%%%%%%%%%%%%%%%%%%%%%%%%%%%%
% EXPERIENCE
%%%%%%%%%%%%%%%%%%%%%%%%%%%%%%%%%%%%%

\section{Experience}

\begin{tabularx}{\textwidth}{Xr}
    \textbf{Greenpoint Global} -- Data Engineering Intern
        & \textit{Jun 2026 -- Jul 2026}
\end{tabularx}

\begin{onecolentry}
\begin{highlights}
    \item Architected Python and SQL ETL pipelines powering analytics across
    \textbf{20M+ financial records}, integrating heterogeneous Apollo, FFIEC,
    and NCUA datasets into standardized, analytics-ready relational models.
    \item Engineered automated data-quality and transformation workflows to
    detect missing, duplicate, and inconsistent records before downstream
    analysis, improving reliability across large-scale financial datasets.
    \item Built dimensional data models, Power BI dashboards, and DAX measures
    that transformed raw banking data into actionable financial, operational,
    and portfolio-level insights.
\end{highlights}
\end{onecolentry}

\vspace{0.10cm}

\begin{tabularx}{\textwidth}{Xr}
    \textbf{Eveoy} -- Software Engineering Intern
        & \textit{Sep 2025 -- Nov 2025}
\end{tabularx}

\begin{onecolentry}
\begin{highlights}
    \item Shipped application features across the stack, connecting frontend
    workflows with Python REST APIs, backend application logic, and SQL-backed
    data services.
    \item Owned cross-stack debugging by tracing user interactions through API
    requests, backend handlers, SQL queries, and response mappings to resolve
    integration and data-flow failures.
    \item Strengthened application reliability with API and SQL validation checks
    that detected inconsistent records and verified data integrity across
    request-response workflows.
\end{highlights}
\end{onecolentry}

%%%%%%%%%%%%%%%%%%%%%%%%%%%%%%%%%%%%%
% PROJECTS
%%%%%%%%%%%%%%%%%%%%%%%%%%%%%%%%%%%%%

\section{Projects}

\begin{tabularx}{\textwidth}{Xr}
    \textbf{AIStora -- Agentic Data Analytics Platform}
        & \href{https://github.com/JamieWei213213/AiStora}{GitHub}\\
    \textit{Python, Flask, PostgreSQL, AWS, Terraform, Docker, GitHub Actions}
        & \\
\end{tabularx}

\begin{onecolentry}
\begin{highlights}
    \item Built a privacy-first analytics product that translates customer-style
    business questions into multi-step plans across 11 typed tools while keeping
    raw financial data out of the LLM context.
    \item Deployed the containerized application on AWS ECS/Fargate behind an
    Application Load Balancer, with encrypted datasets in S3, application data
    in private RDS PostgreSQL, and operational logs in CloudWatch.
    \item Provisioned a two-Availability-Zone VPC with public, private application,
    and isolated database subnets; restricted traffic using security groups and
    an S3 VPC endpoint and separated least-privilege ECS IAM roles.
    \item Automated infrastructure and delivery with Terraform and GitHub Actions
    using AWS OIDC, ECR, rolling ECS deployments, health checks, and rollback;
    validated application behavior with 47 automated tests.
\end{highlights}
\end{onecolentry}

\vspace{0.10cm}

\begin{tabularx}{\textwidth}{Xr}
    \textbf{Cartoonify Image Translation -- CycleGAN}
        & \href{https://github.com/JamieWei213213/Cartoonify}{GitHub}\\
    \textit{Python, PyTorch, Deep Learning, Computer Vision}
        & \\
\end{tabularx}

\begin{onecolentry}
\begin{highlights}
    \item Engineered a CycleGAN in PyTorch for unpaired photo-to-cartoon image
    translation using adversarial, cycle-consistency, and identity objectives.
    \item Built preprocessing and training pipelines for \textbf{20,000+ images}
    and reduced test-set \textbf{FID from 115 to 82.04} through systematic
    loss-weight and learning-rate tuning.
\end{highlights}
\end{onecolentry}

%%%%%%%%%%%%%%%%%%%%%%%%%%%%%%%%%%%%%
% SKILLS
%%%%%%%%%%%%%%%%%%%%%%%%%%%%%%%%%%%%%

\section{Technical Skills}

\begin{onecolentry}
    \textbf{Languages:} Python, SQL, C\#, JavaScript

    \vspace{1pt}

    \textbf{Cloud, Networking \& Security:}
    AWS (ECS/Fargate, S3, RDS, ECR, IAM, VPC, CloudWatch), Terraform,
    Load Balancing, Subnets, Security Groups, Docker, GitHub Actions, CI/CD

    \vspace{1pt}

    \textbf{Backend \& Software Engineering:}
    Flask, SQLAlchemy, REST APIs, PostgreSQL, Linux, Git, Pytest, Debugging

    \vspace{1pt}

    \textbf{Machine Learning \& Agentic AI:}
    Gemini API, LLM Function Calling, Agentic Workflows, Prompt Engineering,
    PyTorch, Scikit-learn

    \vspace{1pt}

    \textbf{Data Engineering \& Analytics:}
    Spark, Pandas, NumPy, ETL, Data Cleaning, Data Modeling,
    Power BI, Tableau
\end{onecolentry}

\end{document}
